% How to cite this draft — appended to the front matter so the citation
% information travels with every PDF copy.
\thispagestyle{empty}
\vspace*{1.5em}
\noindent{\large\bfseries\sffamily How to cite this draft}\\[8pt]

\noindent\small
This is the \textbf{Public Draft v0.9 (Open Review Edition)} of \emph{Machine Learning by Design: From Problem Framing to Reliable Systems}. Each Zenodo deposit also carries a specific \emph{version} DOI; cite the version DOI of the deposit you actually used when reproducibility matters. The \emph{concept} DOI below resolves to the latest version of the project and is the right link for general references.\\[6pt]

\noindent Urmian, S. (2026). \emph{Machine Learning by Design: From Problem Framing to Reliable Systems} (Public Draft v0.9). Zenodo. \url{https://doi.org/10.5281/zenodo.19341954}\\[10pt]

\noindent\textbf{BibTeX --- book draft}
\begin{verbatim}
@book{urmian2026mlbd,
  author    = {Urmian, Sam},
  title     = {Machine Learning by Design: From Problem
               Framing to Reliable Systems},
  year      = {2026},
  publisher = {Zenodo},
  version   = {Public Draft v0.9 (Open Review Edition)},
  doi       = {10.5281/zenodo.19341954},
  url       = {https://github.com/mlgorithm/ml-by-design},
  note      = {Concept DOI shown; cite the specific
               version DOI from the Zenodo deposit when
               reproducibility matters.}
}
\end{verbatim}

\noindent\textbf{BibTeX --- companion code}
\begin{verbatim}
@software{urmian2026mlbd_code,
  author    = {Urmian, Sam},
  title     = {Machine Learning by Design: Companion Code},
  year      = {2026},
  publisher = {Zenodo},
  version   = {0.9.0},
  url       = {https://github.com/mlgorithm/ml-by-design},
  license   = {MIT},
  note      = {Cite the version DOI of the GitHub release
               archived on Zenodo when reproducibility
               matters.}
}
\end{verbatim}

\clearpage
